% ECONOMICS_PROBLEM: {"id": "I11-526", "title": "Cross-market health-care consolidation", "jel_code": "I11", "source_id": "OP-0440"}
% !TeX program = xelatex
\documentclass[11pt,a4paper]{article}
\usepackage[margin=23mm]{geometry}
\usepackage{fontspec}
\setmainfont{Latin Modern Roman}
\usepackage{amsmath,amssymb,mathtools}
\usepackage{xcolor,enumitem,booktabs,longtable,array,xurl,hyperref}
\definecolor{muted}{HTML}{53636C}
\hypersetup{colorlinks=true,urlcolor=blue,linkcolor=blue}
\setlength{\parindent}{0pt}
\setlength{\parskip}{5pt}
\newcommand{\R}{\mathbb R}
\newcommand{\E}{\mathbb E}
\newcommand{\N}{\mathbb N}
\newcommand{\ind}{\mathbf 1}
\newcommand{\Var}{\operatorname{Var}}
\newcommand{\Cov}{\operatorname{Cov}}
\newcommand{\supp}{\operatorname{supp}}
\DeclareMathOperator*{\argmin}{arg\,min}
\DeclareMathOperator*{\argmax}{arg\,max}
\newcommand{\entrymeta}[3]{{\sffamily\small\color{muted}\textbf{\texttt{#1}}\quad|\quad #2\par #3\par}}
\newcommand{\Problemref}[1]{\texttt{#1}}
\newcommand{\JELref}[2]{\texttt{#2} (JEL \texttt{#1})}
\begin{document}
\section*{Cross-market health-care consolidation}
\textbf{Problem ID:} \texttt{I11-526}\quad \textbf{JEL:} \texttt{I11}\par
% BEGIN ECONOMICS STATEMENT
\label{problem:OP-0440}
{\sffamily\small\color{muted}Original subject group: Health economics\par}
\label{definitions:problem:30007722}
\textbf{Audit classification:} Background; proposed extension. RAND identity and broad evidence gaps verified, but the exact p.13 geographic-proximity passage was not recovered; downgraded that claim.
\subsection*{Definitions and precise research target}
\textit{Editorial formalization.} Consider United States hospital mergers between systems operating in different hospital referral regions during 2010–2025. Let \(D_m=1\) denote completion of merger \(m\), and let \(d_m\geq0\) be the premerger distance between the nearest facilities of the merging systems. For commercially insured baseline patients in the affected systems, let \(P_{imh}(D_m)\) be a risk-adjusted negotiated price for a fixed basket of services \(h\) years after the transaction. Define
\[\tau_h(d)=E[\log P_{imh}(1)-\log P_{imh}(0)\mid d_m=d],\quad h\in\{1,3,5\}.\]
Assume the compared prices are positive and log-price contrasts integrable on the stated distance support. The expectation retains the baseline patient mix and service basket, so changing case composition does not mechanically create a price effect. Identify the distance-response profile using a stated merger counterfactual, and distinguish insurer bargaining connections from direct patient substitution. Conditional parallel trends or an alternative quasi-experimental restriction must be justified, with anticipation and contemporaneous acquisitions addressed. Estimate quality and access separately before drawing welfare conclusions. Determine whether any policy-relevant distance range has materially positive price effects, reporting uncertainty rather than declaring a universal geographic cutoff.

Let a merger be the assignment unit, \(I_m\) its baseline patient cohort, and \(P_{imh}(a)\) the price of a basket fixed before the transaction, deflated to a common year. No-merger potential outcomes specify continued independent systems and how other acquisitions evolve. Distance \(d_m\) is measured before announcement. If distance is continuous, the regular conditional expectation \(\tau_h(d)\) is defined only for almost every \(d\) under its distribution; pointwise interpolation requires a continuity or smoothing restriction. A more directly identified target is the interval average \(\tau_h(B_d)=E[\log P(1)-\log P(0)\mid d_m\in B_d]\) for prespecified bins of positive probability. Merger propensity, patient weighting, and joint-merger interference must be declared. Price changes are transfers as well as cost signals and are not by themselves welfare changes. All expectations use a stated baseline population and probability law. Identification means every admissible data-generating model with the same observed-data law yields the same displayed target; otherwise report the identified set. Exact equations, horizons and assumptions are editorial, rather than source-given canonical conjectures.
\subsection*{Variants and assumption changes}
\begin{enumerate}
\item \textbf{Editorial network variant.} Add premerger insurer-network overlap \(n_m\); estimate \(\tau_h(B_d,B_n)\) in positive-probability distance/overlap bins. Geography alone may omit bargaining exposure.
\item \textbf{Editorial welfare variant.} Replace prices with net-health benefit \(\lambda\Delta Q-\Delta C_{\rm resource}-\rho\Delta T\). A positive price effect does not establish this contrast's sign.
\end{enumerate}
\subsection*{References and source audit}
\textbf{1.} Consolidation quality and access evidence gaps verified; claimed exact geographic-proximity passage not independently retrieved. Locator: RAND RR-A1820-1, 30 September 2022; report-page Key Takeaways. Report metadata and primary RAND summary inspected; PDF follow-up retrieval timed out. The supplied p.13 proximity locator remains unverified. URL: \url{https://www.rand.org/content/dam/rand/pubs/research_reports/RRA1800/RRA1820-1/RAND_RRA1820-1.pdf}.
\begin{enumerate}\raggedright
\item\hypertarget{definitions:source:30007722:1}{} Jodi L. Liu; Zachary M. Levinson; Annetta Zhou; Xiaoxi Zhao; PhuongGiang Nguyen; Nabeel Qureshi. 2022. \emph{Environmental Scan on Consolidation Trends and Impacts in Health Care Markets}. RAND RR-A1820-1, 30 September 2022; report-page Key Takeaways.
\url{https://www.rand.org/content/dam/rand/pubs/research_reports/RRA1800/RRA1820-1/RAND_RRA1820-1.pdf}.
DOI: \url{https://doi.org/10.7249/RRA1820-1}.
\end{enumerate}

\subsection*{Common conventions: Source mathematical conventions}

These conventions apply to each entry unless explicitly replaced.

\textbf{Models and identification.} A primitive model \(m\) specifies all probability laws, feasible choices, information, equilibrium and measurement rules used by its target. The admissible class \(\mathcal M\) is fixed independently of outcomes. If \(P_O(m)\) is its observable probability law and \(\tau(m)\) the target, its identified set at observable law \(P\) is
\[
 \mathcal I_\tau(P)=\{\tau(m):m\in\mathcal M,\ P_O(m)=P\}.
\]
Point identification means that this set is a singleton. An empty set signals incompatibility of data and assumptions. Coordinate bounds need not describe the joint set. Bounds are sharp only if every retained target is attainable or arbitrarily approachable by compatible models. Finite bounds require explicit restrictions on unobserved quantities; unrestricted confounding, hidden wealth, extrapolation or arbitrary equilibrium selection can yield uninformative sets.

\textbf{Causal interventions.} A potential outcome \(Y_i(p)\) is the outcome under a complete policy regime \(p\), including timing, financing, information and market spillovers. The target population and units are fixed before treatment. With interference, use the complete assignment vector or a justified exposure mapping; individual no-interference notation alone cannot identify an economy-wide intervention. A difference of mean potential outcomes does not require the cross-world joint distribution. Conditioning on post-treatment migration, survival, enrollment or employment changes the estimand and needs separate justification.

\textbf{Statistical statements.} An estimator, confidence set, forecast or test specifies a sampling law, model class and loss/coverage criterion. Uniform \((1-\alpha)\)-coverage means \(\inf_{m\in\mathcal M}\Pr_m\{\tau(m)\in C_n\}\geq1-\alpha\), or an explicitly asymptotic version. The whole parameter space trivially has coverage, so informativeness requires a stated width, risk or power objective. A forecasting improvement is relative to a named benchmark, information set and evaluation population.

\textbf{Equilibrium and optimization.} An equilibrium comprises feasible plans, optimality, beliefs where needed, budget constraints and all market/resource clearing. The primitives or a stated benchmark define each correspondence \(\mathcal E(p;m)\). Nonemptiness is assumed or proved, never inferred just from compact policy sets. A selected-equilibrium target uses a declared selection rule; a robust target quantifies over all permitted equilibria. Use suprema or infima unless attainment is established. An \(\varepsilon\)-optimal policy is feasible and within \(\varepsilon\) of the finite optimum value. Report an empty feasible set separately.

\textbf{Welfare and accounting.} Normative utility, interpersonal normalization, social weights, numeraire, discounting and the use of public receipts are primitives. Behavioral utility need not coincide with the welfare criterion. Household welfare includes ownership income and rents once; adding profits again to compensating variation generally double-counts them. A monetary equivalent is defined by an explicit transfer and price convention, with monotonicity and range conditions. Average effects, mechanism contributions, transfers and welfare losses are different targets. Infinite sums require convergence and dynamic feasible plans require terminal or no-Ponzi conditions.

\textbf{Source versions.} A working-paper date, revision date and journal publication year are distinguished. A DOI for a journal version is not automatically the DOI of an earlier manuscript. Locators refer to the stated version and printed pages where specified. A metadata-only check does not verify a claim about a full-text conclusion. Every entry's source subsection narrows the provenance claim accordingly.
% END ECONOMICS STATEMENT
\end{document}
